\section{First variations of equality and the original-data adjoint}
\label{eq-section}

In this section \(S\) is connected, is an outermost future MOTS, and is
outer area-minimizing in the full-cut sense.  We assume equality in
Theorem~\ref{thm:ah-invariant}.  All variations below are variations of the
original pair \((g,K)\) on the original manifold \(\Omega\).  Nearby data
need satisfy only the hypotheses of the numerical inequality; their
boundaries need not remain outermost or area-minimizing.
Here \(\tau\) again denotes the original asymptotically hyperbolic decay
exponent in Definition~\ref{def:setup-data}.

Choose a balanced spatial chart, so that \(p_0=m_{\rm AH}\) and \(p_i=0\)
at the data under consideration.  Put
\begin{equation}
 A=\Area_g(S),\qquad r_h=\sqrt{A/(4\pi)},\qquad
 c=16\pi\frac{d}{dA}\frac{r_h+r_h^3}{2}
   =\frac{1+3r_h^2}{r_h},\qquad \kappa_h=\frac c2 .
 \label{eq-area-constant}
\end{equation}
The fixed-chart time component is a supporting functional for the
Lorentz norm of a future-timelike covector.  Consequently the numerical
inequality gives, for all sufficiently small admissible variations,
\begin{equation}
 p_0(g_s)\ \geq\ m_{\rm AH}(g_s)\
 \geq\ \frac{r_{A(s)}+r_{A(s)}^3}{2},
 \qquad r_{A(s)}=\sqrt{A(s)/(4\pi)}.
 \label{eq-supporting-inequality}
\end{equation}
Here \(A(s)\) is the minimum full enclosing area for \(g_s\).  Equality
holds throughout at \(s=0\).  Only spatial chart covariance is used in
choosing the balanced chart.

\begin{theorem}[Causal Killing data produced by equality]
\label{eq-causal-killing}
There are a smooth function \(u\) and a smooth vector field \(X\) on the
one-sided manifold \(\Omega\) with the following properties:
\begin{enumerate}
\item \(u\geq |X|_g\) everywhere and \(u>0\) in \(\operatorname{int}\Omega\).
For some \(3/2<\beta<\min\{2,\tau\}\),
\begin{equation}
 u-V_0=O_{2,\alpha'}(r^{1-\beta}),\qquad
 X=O_{2,\alpha'}(r^{1-\beta})
 \label{eq-end-normalization}
\end{equation}
in background tensor norm, with \(0<\alpha'<\alpha\).
\item In the interior, and by continuity at the boundary,
\begin{align}
 \operatorname{sym}\nabla X&=-uK, \label{eq-shift}\\
 \Delta u-3u&=-\operatorname{div}_g(K(X,\cdot)^\sharp),\label{eq-lapse}\\
 \Hess_g u&=
 u(\Ric_g+3g+(\tr_gK)K-2K^2)
 -\mathcal L_XK+8\pi X_{(i}J_{j)}. \label{eq-full-lapse}
\end{align}
Symmetrization includes averaging.  Moreover \(u\mu+J(X)=0\).
\item On \(\mathbb R\times\operatorname{int}\Omega\), the metric
\begin{equation}
 \mathbf g=-u^2dt^2+
 g_{ij}(dx^i+X^i\,dt)(dx^j+X^j\,dt)
 \label{eq-stationary-metric}
\end{equation}
induces precisely the original \(g\) and the original \(K\) on \(t=0\),
with future normal \(n=u^{-1}(\partial_t-X)\).  Writing
\(\xi=\partial_t\) and \(N=u^2-|X|_g^2\), one has
\begin{equation}
 \Ric_{\mathbf g}=-3\mathbf g+
       8\pi\mu u^{-2}\xi^\flat\otimes\xi^\flat,\qquad
 \mu N=0 .
 \label{eq-null-matter}
\end{equation}
\item For the normal \(\nu\) pointing into the exterior,
\begin{equation}
 X=u\nu,\qquad
 \partial_\nu u+K(X,\nu)=\kappa_h\quad\hbox{on }S.
 \label{eq-horizon-data}
\end{equation}
Either \(u>0\) on all of \(S\), or \(u=0\) on all of \(S\).
\end{enumerate}
\end{theorem}

We prove the theorem without assuming a stationary development.  The
Lorentzian metric in its statement will be constructed only after the
multiplier fields have been obtained and shown to be smooth and positive.

\subsection{The envelope derivative for full enclosing area}

Extend \(g\) smoothly across a short collar of \(S\), and fix the entire
inner side of that collar as an obstacle.  This auxiliary extension is
used only for perimeter minimization.  It is not required to satisfy
constraints.  Minimize among bounded filled sets containing the
obstacle, taking perimeter in the extended open collar and the original
exterior.  Thus the obstacle itself has boundary \(S\) and perimeter
\(\Area_g(S)\): the convention does not erase \(S\) when the exterior
competitor is \(D=\Omega\).

Bounded complementary pockets can be filled.  Their removal from the
exterior component cannot increase perimeter, and any nonempty pocket
of a regular minimizer would contribute removable positive perimeter.
The remaining complement is the connected exterior containing the end.
Conversely, every full cut in the original definition gives a filled
competitor in this formulation.  A compact \(C^{1,1}\) obstacle boundary
can be pushed slightly into the exterior by a collar vector field and
then smoothed in graph charts.  These operations preserve separation
and make its areas converge.  The perimeter infimum therefore equals
the infimum over the smooth full cuts of Definition~\ref{def:setup-cut}.

Let \(\mathcal M\) denote the family of boundaries attaining this
perimeter infimum for \(g\).

\begin{lemma}[Compact minimizing family and area derivative]
\label{eq-area-envelope}
Let \(g_s\) be a smooth metric path through \(g\), with its first two
parameter derivatives bounded relative to \(g\), and with a common
asymptotically hyperbolic \(C^{2,\alpha'}\) bound of positive decay.
Then, after restricting the parameter interval, the minimizing
boundaries for \(g_s\) lie in a common compact set.  The family
\(\mathcal M\) is compact for indicator \(L^1\) convergence and
tangent-plane area-measure convergence.  If \(H=\dot g_0\), then
\begin{equation}
 \left.\frac{d}{ds}\right|_{0+}A(s)
   =\min_{T\in\mathcal M}a'_T(H),\qquad
 a'_T(H)=\frac12\int_T\tr_T H\,dA_g.
 \label{eq-area-derivative}
\end{equation}
Off the obstacle, minimizing boundaries that meet have the same tangent
plane.  This common plane is continuous on their union locally away
from \(S\).
\end{lemma}

\begin{proof}
Large coordinate spheres have mean curvature \(2+o(1)\), uniformly for
these paths.  Fix a radius beyond which this mean curvature is positive
and the obstacle is absent.  On a larger compact truncation, the
direct method for perimeter gives an obstacle minimizer.  Its boundary
cannot have a maximal radius in the strictly mean-convex region: at a
free maximum it is minimal and lies on the inner side of a strictly
mean-convex coordinate sphere, in contradiction with the local
comparison principle.  Contact with the outer truncation can likewise
be moved inward to lower area.  Thus every such minimizer lies inside
the same fixed radius.  Truncations containing any prescribed smooth
competitor show that their minimum is the full infimum.

The local obstacle regularity theorem
\cite[Regularity Theorem~1.3(iii)]{huiskenilmanen2001} applies here to
pure perimeter, a smooth obstacle, a smooth metric and ambient
dimension three.  It gives \(C^{1,1}\) boundaries and smooth minimal
free parts. Part (iii) is used for this individual regularity.
For a fixed \(0<\alpha_0\le1/2\), the local \(C^{1,\alpha_0}\)
estimates in Part (ii) are uniform: the obstacle bounds and positive
metric \(C^1\) bounds are uniform for the family \(g_s\).

Only away from \(S\) do we need density estimates. Here is a local
proof with uniform constants. In a sufficiently small ball centered
at a boundary point, compare a minimizing set \(E\) with
\(E\setminus B_\rho\) and \(E\cup B_\rho\), filling any unnecessary
pockets afterwards. The ball misses the obstacle. For
\(f(\rho)=\Vol(E\cap B_\rho)\), minimality and coarea give, for
almost every \(\rho\),
\[
 P(E;B_\rho)\le f'(\rho),\qquad
 P(E\cap B_\rho)\le2f'(\rho).
\]
Uniform local isoperimetry implies
\(f'(\rho)\ge c f(\rho)^{2/3}\), hence
\(f(\rho)\ge c\rho^3\). The same argument applies to the complement.
Relative isoperimetry then gives
\(P(E;B_\rho)\ge c\rho^2\); the comparison above and the uniform
sphere-area bound give \(P(E;B_\rho)\le C\rho^2\).
These constants are uniform on compact sets disjoint from \(S\).
Perimeter compactness now gives subsequential indicator convergence.  Lower
semicontinuity, together with testing a fixed minimizer, gives
convergence of perimeters, including for varying metrics.

Here is the measure-continuity detail needed for differentiation.
In a fixed orthonormal frame let \(D\chi_{E_j}\) be the vector normal
measures.  Their weak convergence and convergence of total variations
imply weak convergence of their variation measures: any weak limit
dominates \(|D\chi_E|\), and equality of total masses eliminates the
excess.  Approximate the limiting unit normal in the square-integral
norm of this measure by continuous vector fields.  For such a field
\(v\), expand
\[
 \int|\nu_j-v|^2\,d|D\chi_{E_j}|
 =\int(1+|v|^2)\,d|D\chi_{E_j}|
       -2\int v\cdot dD\chi_{E_j}.
\]
Both terms converge.  Continuous approximation and uniform continuity
then give convergence for every continuous function of position and
unit normal, and hence for functions of the unoriented tangent plane.
A partition of unity proves the assertion globally.  Uniform metric
convergence converts varying-metric perimeters to this fixed-metric
argument.

On every rectifiable cut of bounded area, the area-element expansion is
uniform:
\begin{equation}
 \Area_{g_s}(T)=\Area_g(T)+s\,a'_T(H)
                      +O(s^2)\Area_g(T).
 \label{eq-area-taylor}
\end{equation}
A fixed minimizing cut gives the upper one-sided bound in
\eqref{eq-area-derivative}.  For the lower bound take minimizers \(T_s\)
for \(g_s\), pass along any sequence \(s\downarrow0\) to a member
\(T\in\mathcal M\), and use
\[
 \frac{A(s)-A(0)}s
   \geq a'_{T_s}(H)+O(s)\longrightarrow a'_T(H).
\]
Compactness and continuity of \(a'_T(H)\) prove the formula.

Finally, the union and intersection of two filled minimizing sets
remain obstacle competitors.  Perimeter submodularity gives
\[
 P(E_1\cup E_2)+P(E_1\cap E_2)
       \leq P(E_1)+P(E_2)=2A.
\]
Both terms on the left are at least \(A\), so both are minimizers.
Different tangent planes at an intersection would give a corner in
one of them, contradicting interior regularity.  To check continuity
on the union, take points on a sequence of minimizers converging away
from the obstacle.  The two phase-volume density bounds prevent disappearance of the
limiting point under indicator convergence, so it lies on a limiting
minimizer. The uniform local \(C^{1,\alpha_0}\) estimates and strict
perimeter convergence then give tangent-plane convergence there.  The
common-plane property makes this independent of the chosen sequence.
\end{proof}

\subsection{A strict direction and the asymptotic test generators}

Use the fixed unit ball bundle
\[
 \mathcal B=\{(x,v):x\in\Omega,\ |v|_g\leq1\},\qquad
 \mathcal C(x,v)=16\pi(\mu(x)+J_x(v)),
 \qquad \mathcal Z=\{\mathcal C=0\}.
\]
In a metric variation \(g_s(\cdot,\cdot)=g(A_s\cdot,\cdot)\), transport
\(v\) by \(v_s=A_s^{-1/2}v\).  This preserves its length and gives
\begin{equation}
 \dot v_0=-\tfrac12 H^\sharp v.
 \label{eq-vector-transport}
\end{equation}
Every constraint derivative below includes this term.

\begin{lemma}[Strict conformal direction]
\label{eq-strict-direction}
Choose \(3/2<\beta<\min\{2,\tau\}\) and
\(0<\delta<\min\{1,2\beta-3\}\).  There are smooth positive weights
\(w,w_2,D\), with \(w=r^{-3-\delta}\), \(w_2=r^{-\beta}\) on a far
end and \(D=O(r^{-\beta})\), \(D\geq |K|\), and a smooth function
\(\phi=O_{2,\alpha'}(r^{-\beta})\), such that
\begin{align}
 (-\Delta+3)\phi
   &=D\sqrt{w_2^2+|d\phi|^2}+w,&
 \partial_\nu\phi&=-1. \label{eq-strict-solve}
\end{align}
It has finite fluxes against all four static potentials.  In
particular,
\[
 (-\Delta+3)\phi\geq |K||d\phi|+w,\qquad
 \frac{(-\Delta+3)\phi}{w}\longrightarrow1.
\]
For \(Q=(4\phi g,2\phi K)\), on the active set one has
\begin{equation}
 \mathcal C'_Q=8\big((3-\Delta)\phi+K(v,\nabla\phi)\big),\quad
 \mathcal C'_Q/w\longrightarrow8,\quad
 \theta'_Q=-4\ \hbox{on }S.
 \label{eq-strict-tests}
\end{equation}
\end{lemma}

\begin{proof}
Solve \eqref{eq-strict-solve} on spherical truncations with outer
Dirichlet value zero.  Continue both the source and the inner Neumann
value from zero.  The gradient term can be written \(b\cdot d\phi+f\),
where
\[
 |b|\leq D,\qquad 0\leq f\leq Dw_2+w .
\]
The operator is consequently \(-\Delta-b\cdot\nabla+3\).
A fixed collar function supplies the prescribed outward Neumann
derivative, which is \(+1\); adding a sufficiently large constant
dominates its bounded compactly supported derivatives and the source.
The maximum principle gives a bound independent of the truncation and
continuation parameter.  Beyond a fixed sphere,
\[
 (3-\Delta_b)r^{-\beta}
   =(3-\beta)(1+\beta)r^{-\beta}+O(r^{-\beta-2}).
\]
The decaying drift and metric errors preserve its positive leading
coefficient.  Comparison with a multiple of this function, using the
already bounded value on the fixed sphere and zero outer data, gives
\(|\phi|\leq C r^{-\beta}\).

Lemma~\ref{lem:linear-separated-faces}, followed by Sobolev embedding
and the one-sided Schauder estimate, gives the same weighted
bound through two derivatives.  At the fixed inner face use the
corresponding Neumann estimates.  The linearized equation has bounded
drift, zeroth-order coefficient \(3\), homogeneous Neumann data at the
inner face and homogeneous Dirichlet data at the outer face.  Its
kernel is zero by the maximum principle, so the usual elliptic
Fredholm alternative makes it invertible.  The a priori estimates
close the method of continuity on each truncation.  Exhaustion and
local compactness then produce the asserted decaying solution, smooth
on compact sets including the one-sided boundary.

The additional source is \(O(r^{-2\beta})=o(w)\), proving the ratio
claim.  Replacing \(\Delta_g\) by \(\Delta_b\) costs
\(O(r^{-\tau-\beta})\).  Both this error and \(r^{-2\beta}\) have
finite integrals with weight \(V_0\,dV_b\), as does \(w\).
The scalar linearization identity
\[
 \operatorname{div}_b\mathbb U(V,4\phi b)
       =8V(3-\Delta_b)\phi,\qquad \Hess_bV=Vb,
\]
therefore proves existence of all the conformal fluxes.
Replacing \(b\) by \(g\) in \(4\phi b\) costs order
\(r^{-\tau-\beta}\), beyond the mass order.

For the exact conformal transformation \(g_s=e^{4s\phi}g\),
\(K_s=e^{2s\phi}K\), direct differentiation gives
\[
 \mathcal C'_Q=-4\phi\mathcal C+
       8((3-\Delta)\phi+K(v,\nabla\phi)).
\]
The first term vanishes on \(\mathcal Z\).
The expansion transforms as
\(\theta_s=e^{-2s\phi}(\theta+4s\partial_\nu\phi)\);
since \(\theta=0\) on \(S\), its derivative is \(-4\).
\end{proof}

In addition to \(Q\), take all compact variations, including arbitrary
one-sided jets at \(S\), and the following end variations, cut off near
the compact interior:
\begin{enumerate}
\item \((4\psi g,2\psi K)\), where
\(\psi=r^{-3}f(\omega)\) and \(f\in C^\infty(\mathbb S^2)\);
\item pure tensor variations \((0,P)\) whose trace reversal with
respect to \(b\) has only mixed block
\[
 (P-(\tr_bP)b)_{\nu_b A}=r^{-3}W_A(\omega)
\]
in radial and spherical orthonormal scaling, for arbitrary smooth
one-forms \(W\) on \(\mathbb S^2\).
\end{enumerate}
Their constraint derivatives are \(o(w)\) on the active rays.
Indeed the conformal radial indicial term cancels, leaving an
\(O(r^{-5})\) background scalar term and products with the decaying
original coefficients.  For the mixed block, its radial derivative
is \(-3(V_0/r)r^{-3}W_A\); the sphere connection terms give
\(+3(V_0/r)r^{-3}W_A\).  Thus its tangential background divergence is
zero, while its normal divergence is \(r^{-4}\operatorname{div}_\sigma W\).
Metric and tensor errors cost \(O(r^{-3-\tau})\).
The choice \(\delta<1\) proves the assertion.

Let \(\mathcal V\) be the real linear space generated by these tests.
The coefficient \(a(H,P)\) of \(Q\) is well defined even if
\(\mathcal Z\) is compact. Indeed, on the full unit ball bundle,
\(\mathcal C=O(r^{-\tau})\), uniformly in \(v\), and the conformal
identity in the preceding proof gives
\[
 \frac{\mathcal C'_Q}{w}
 =8+o(1)-\frac{4\phi\mathcal C}{w}=8+o(1).
\]
Here \(\tau+\beta>2\beta>3+\delta\), so the last term tends to
zero. The other generators have normalized full-ball derivative
 tending to zero by the same end calculations (and compact variations
vanish there). Applying these limits to any linear relation among
the generators forces the coefficient of \(Q\) to vanish. The
active-ray statement is its restriction when active rays exist.

\subsection{Positive separation of the active constraints}

\begin{lemma}[Multiplier identity]
\label{eq-separation}
There are a probability measure \(\omega\) on \(\mathcal M\), a
nonnegative finite measure \(\eta\) on \(S\), a nonnegative locally
finite measure \(\mathfrak L\) on \(\mathcal Z\), and \(z\geq0\), such
that \(w\mathfrak L\) is finite and
\begin{equation}
 16\pi p_0'-c\int_{\mathcal M}a'_T(H)\,d\omega(T)
 =\langle\mathcal C',\mathfrak L\rangle
       -\langle\theta',\eta\rangle+z\,a(H,P)
 \label{eq-multiplier-identity}
\end{equation}
for every \((H,P)\in\mathcal V\).
\end{lemma}

\begin{proof}
Compactify \(\mathcal Z\) by one point, collapsing all rays whose base
points escape compact sets; if it was compact, add an isolated point.
The ball fibers are compact.  The functions \(\mathcal C'/w\) extend
continuously with value \(8a(H,P)\).  On the disjoint compact union of
this space, \(S\), and \(\mathcal M\), consider the linear test map
\begin{equation}
 (H,P)\longmapsto
 \big(\mathcal C'/w,\ -\theta',\
          c\,a'_T(H)-16\pi p_0'\big).
 \label{eq-test-map}
\end{equation}
Lemma~\ref{eq-area-envelope} supplies continuity on the last piece.

Suppose its image contained a function positive everywhere.  Write its
direction as \(aQ+(H_0,P_0)\).  Positivity at the additional point gives
\(a>0\).  Realize the direction by
\[
 g_s=e^{4as\phi}(g+sH_0),\qquad
 K_s=e^{2as\phi}(K+sP_0).
\]
These metrics remain positive and uniformly comparable to \(g\).
On the far end, retain the nonnegative background constraint with its
positive conformal scaling factor.  The remaining terms have the
following bounds, uniformly over the unit ball:
\[
 \begin{array}{c|c}
 \text{term}&\text{size}\\ \hline
 \text{strict conformal gain}&8asw\\
 \text{non-strict linear changes}&s\,o(w)\\
 \text{new quadratic and cross terms}&O(s^2w).
 \end{array}
\]
For clarity, the non-strict end metric and tensor tests are
\(O(r^{-3})\), their background constraint terms are
\(O(r^{-4})+O(r^{-3-\tau})\), and their products with the strict
direction are \(O(r^{-3-\beta})\).  The strict direction's own
quadratic errors are \(O(r^{-2\beta})\).  All are controlled by
\(w\).  Applying the exact conformal law after the intermediate
variation gives these bounds without differentiating the
nonnegative original constraint into an uncontrolled error.
One may compose the two isometries of unit balls in this calculation:
at an active unit ray a difference of isometric identifications is
tangent to the unit sphere and is annihilated by \(J\); at a vacuum
ray \(J=0\).

On the remaining compact ball bundle, positivity of the derivative
on the closed active set persists on a neighborhood of it.  The
original constraint has a positive minimum on the compact complement.
Taylor's formula thus gives the DEC for sufficiently small \(s>0\)
everywhere.  The boundary has strictly negative future expansion.
The same estimates are integrable with weight \(V_0\), and the flux
divergence identity proves differentiability of all four fluxes.
The covector remains future timelike by continuity.  Thus
\eqref{eq-supporting-inequality} applies to this path.

Positivity on \(\mathcal M\), however, says
\[
 c\min_{T\in\mathcal M}a'_T(H)-16\pi p'_0>0.
\]
By Lemma~\ref{eq-area-envelope} and \eqref{eq-area-constant}, the
right derivative at zero of the difference in
\eqref{eq-supporting-inequality} is negative.  That difference is
initially zero and is nonnegative for \(s>0\), a contradiction.

The image of \eqref{eq-test-map} is therefore disjoint from the open
positive cone of continuous functions.  Hahn--Banach separation
gives a nonzero continuous functional, nonnegative on nonnegative
functions and annihilating this linear image.  The Riesz
Representation Theorem gives finite nonnegative measures on its
three compact pieces.  The objective piece has positive mass:
otherwise evaluation on \(Q\), strictly positive on both other
pieces, would contradict annihilation.  Normalize that mass to one.
Divide the finite measure on finite active rays by \(w\) and call
it \(\mathfrak L\); absorb eight times the end-point atom into \(z\).
Rearranging the identity gives \eqref{eq-multiplier-identity}.
\end{proof}

Define the first moments, initially as measures with the fixed-volume
distribution convention, by
\[
 \langle u,f\rangle=\int f(x)\,d\mathfrak L(x,v),\qquad
 \langle X,\alpha\rangle=\int\alpha_x(v)\,d\mathfrak L(x,v).
\]
Positivity and activity give
\begin{equation}
 u\geq |X|_g,\qquad u\mu+J(X)=0
 \label{eq-measure-complementarity}
\end{equation}
as measure identities and inequalities.

\subsection{Removing the interior area multipliers}

\begin{lemma}[Normal slice tests]
\label{eq-normal-tests}
For every \(s\in C_c^\infty(\operatorname{int}\Omega)\), there are
compact variations \((H_\epsilon,P_\epsilon)\) with
\(H_\epsilon=2sK\), supported in one fixed compact set, whose
active-ray constraint derivatives tend uniformly to zero.
\end{lemma}

\begin{proof}
On a compact neighborhood of the support, choose a Gaussian
Lorentzian collar
\[
 \mathbf g_\epsilon=-dt^2+g_\epsilon(t),\qquad
 g_\epsilon(0)=g,\quad \partial_tg_\epsilon(0)=2K.
\]
Prescribe its second time jets so that
\[
 \frac1{8\pi}(G_{\mathbf g_\epsilon}-3\mathbf g_\epsilon)_{ij}
       =D_{\epsilon,ij},\qquad
 D_\epsilon=\frac{J\otimes J}{\mu+\epsilon}
 \quad\hbox{at }t=0.
\]
The spatial Einstein tensor depends on the freely chosen normal
derivative of \(K\) by a nonzero multiple of trace reversal.  In
dimension three trace reversal has eigenvalues \(1\) and \(-2\);
it is invertible.  The remaining terms are fixed spatial jets.
Thus these smooth second jets can be prescribed, and a quadratic
time extension with a cutoff realizes a smooth Lorentzian metric
on a sufficiently short collar.  The normal and mixed components
of the displayed tensor are \(\mu,J\), by the constraints.

Set \(D=J\otimes J/\mu\) where \(\mu>0\), and set it to zero at
vacuum.  The DEC gives
\[
 |D_\epsilon|\leq\mu,\qquad
 |\nabla D_\epsilon|\leq2|\nabla J|+|\nabla\mu|.
\]
At a zero of the smooth nonnegative function \(\mu\), its differential
vanishes; \(|J_i|\leq\mu\) then gives \(dJ_i=0\) there also.
The displayed derivative bound proves that \(D\) is \(C^1\) across
vacuum with zero derivative there.  Splitting any compact set into
\(\{\mu\geq\delta\}\) and \(\{\mu<\delta\}\) proves
\(D_\epsilon\to D\) in \(C^1\).

Vary the slice by graphs \(t=a s(x)\), differentiating at \(a=0\).
The metric derivative is \(2sK\).  At a nonvacuum active ray,
\(\mu=|J|\), \(v=-J^\sharp/\mu\), and \(\zeta=n+v\) is null.
The limiting prescribed stress has the form
\[
 \mu^{-1}\gamma\otimes\gamma,\qquad
 \gamma(n)=\mu,\quad\gamma|_{T\Omega}=J,\quad\gamma(\zeta)=0.
\]
The covector \(\gamma\) is nonnegative on future causal vectors.
Parallel transport \(\zeta\) in any spatial direction using the
spacetime connection.  Its contraction with \(\gamma\) is
nonnegative and has a zero, so its derivative there is zero.
Consequently the spatial covariant derivatives of the limiting
stress, contracted in one slot with \(\zeta\), vanish.

Let \(T_\epsilon=(G_{\mathbf g_\epsilon}-3\mathbf g_\epsilon)/(8\pi)\).
Contracted Bianchi, before taking the limit, gives
\[
 (\nabla_nT_\epsilon)(n,\zeta)
       =\sum_i(\nabla_{e_i}T_\epsilon)(e_i,\zeta).
\]
The right side tends to zero by the compact \(C^1\) convergence just
proved.  Thus the required normal derivative of the constraint
combination tends to zero too.  Differentiating its arguments causes
no additional limiting term: the first argument is paired against
a tensor with \(\zeta\) in its kernel, and the second argument remains
null; at activity \(\gamma\) is proportional to \(\zeta^\flat\), so
nullity kills its differentiated contraction.  At vacuum the tensor
and its spatial derivatives vanish, and Bianchi gives the same
conclusion for every vector in the closed unit ball.  All expressions
use only the fixed compact jets and \(D_\epsilon,\nabla D_\epsilon\);
the convergence is uniform, including near vacuum.
\end{proof}

\begin{lemma}[Concentration on the original horizon]
\label{eq-area-concentration}
The measure \(\omega\) in Lemma~\ref{eq-separation} is concentrated on
the single cut \(S\).
\end{lemma}

\begin{proof}
Apply Lemma~\ref{eq-normal-tests} in
\eqref{eq-multiplier-identity}.  Mass, expansion and end-point terms
are zero.  The constraint measure is finite over the fixed compact
support, so the limit gives
\[
 \int_{\mathcal M}\int_T s\,\tr_TK\,dA_g\,d\omega(T)=0
       \qquad(s\in C_c^\infty(\operatorname{int}\Omega)).
\]
Off \(S\), form the positive aggregate area measure
\(\mathfrak b=\int_{\mathcal M}dA_T\,d\omega(T)\).
Lemma~\ref{eq-area-envelope} makes
\(F(x)=\tr_{\Pi(x)}K(x)\), for the common tangent plane \(\Pi(x)\),
a single continuous function on the union of the cuts locally
away from \(S\).  The last identity is \(F\mathfrak b=0\).
Fubini on a countable compact exhaustion shows that, for
\(\omega\)-almost every cut, \(\tr_TK=0\) almost everywhere off
\(S\).  Such a cut is smooth minimal there, so continuity gives
\(H_T=\tr_TK=0\) throughout its free part.

At contact, write the cut and \(S\) as \(C^{1,1}\) graphs with the
same orientation.  Their second derivatives agree almost everywhere
on the coincidence set: apply twice the fact that the derivative
of a Lipschitz function vanishing on a measurable set is zero
almost everywhere on that set.  The cut therefore solves
\(H_T+\tr_TK=0\) almost everywhere across contact as well as away
from it.  This is a uniformly elliptic quasilinear graph equation.
Its coefficients are Hölder after the \(C^{1,1}\) bound; interior
elliptic regularity and differentiation make the graph smooth.

If this smooth MOTS touches \(S\), their one-sided difference
satisfies a linear elliptic equation with bounded coefficients.
The strong comparison principle gives local coincidence.
The contact set is thus open and closed in connected \(S\), so
the cut contains \(S\).  Its full area is already \(A\), excluding
additional positive-area components.  If it does not touch \(S\),
it is a smooth compact full enclosing MOTS entirely in the interior,
after filling removable pockets.  This contradicts outermostness.
Hence almost every cut is exactly \(S\).
\end{proof}

\subsection{A local adjoint lemma and finite-type continuation}

The following version records the cosmological constant explicitly,
for later use on annuli.  Set
\[
 16\pi\mu_{\Lambda_{\rm c}}
    =R-2\Lambda_{\rm c}+(\tr K)^2-|K|^2,\qquad
 8\pi J=\operatorname{div}_g(K-(\tr K)g).
\]

\begin{lemma}[Local regularity and the full adjoint]
\label{eq-local-adjoint}
Let \(g,K\) be smooth on a connected open three-manifold and let
\(\Lambda_{\rm c}\) be constant.  Suppose locally finite moment
measures \(u,X\) satisfy \(u\geq|X|\),
\(u\mu_{\Lambda_{\rm c}}+J(X)=0\), and annihilate every compact
variation of \(16\pi(\mu_{\Lambda_{\rm c}}+J(v))\), including the
isometric-vector correction \eqref{eq-vector-transport}.
Then \(u,X\) are smooth and
\begin{align}
 \operatorname{sym}\nabla X&=-uK,\label{eq-general-shift}\\
 \Delta u+\Lambda_{\rm c}u
   &=-\operatorname{div}_g(K(X,\cdot)^\sharp),\label{eq-general-lapse}\\
 \Hess u
   &=u(\Ric-\Lambda_{\rm c}g+\tau_KK-2K^2)
       -\mathcal L_XK+8\pi X_{(i}J_{j)},\qquad \tau_K=\tr K.
       \label{eq-general-full-lapse}
\end{align}
Their full first jet satisfies a homogeneous linear first-order
system with smooth coefficients depending on \(g,K,\Lambda_{\rm c}\).
A jet vanishing at one point vanishes everywhere.  If the moments
are nonzero, then \(u>0\) everywhere.
\end{lemma}

\begin{proof}
For a pure tensor variation \(P\), integration of momentum by parts
gives the distributional coefficient
\[
 2\big[u(\tau_Kg-K)-\operatorname{sym}\nabla X
                         +(\operatorname{div}X)g\big].
\]
Its trace gives \(\operatorname{div}X=-u\tau_K\), and substitution
gives \eqref{eq-general-shift}.  A simultaneous conformal test
\((4\psi g,2\psi K)\), using complementarity to cancel its background
scaling term, gives
\[
 8\langle u,(-\Delta-\Lambda_{\rm c})\psi\rangle
       +8\langle X,K(\nabla\psi,\cdot)\rangle=0.
\]
This is \eqref{eq-general-lapse}.  Diverging the shift equation and
using its trace gives
\begin{equation}
 \Delta X_j+\Ric_{jk}X^k
  =(\tau_K\delta_j{}^i-2K_j{}^i)\nabla_i u
       +u(\nabla_j\tau_K-2\nabla^iK_{ij}).
 \label{eq-vector-elliptic}
\end{equation}
Together with the lapse equation this is an elliptic system with
diagonal scalar-Laplacian principal part and smooth first-order
couplings.  Measures belong locally to a negative Sobolev space.
The local elliptic estimate on nested compact cutoffs improves
their Sobolev order by one, because the couplings have at most
one derivative.  Iteration and Sobolev embedding yield smoothness.

For completeness, the metric equation is computed next.  Put
\(P_K=K-\tau_Kg\).  The compact constraint pairing is the variation of
\[
 \int\big[u(R-2\Lambda_{\rm c}+\tau_K^2-|K|^2)
                  +2X^i\nabla^j(P_K)_{ij}\big]\,dV_g
\]
together with
\(-8\pi\int H(X,J^\sharp)\,dV_g\).
The latter is exactly the vector-transport correction.
Changing the fixed background volume density to the varying one
adds zero, since its background integrand is
\(16\pi(u\mu_{\Lambda_{\rm c}}+J(X))=0\).
Integrating the shift term once gives
\(-\int P_K^{ij}(\mathcal L_Xg)_{ij}\,dV_g\).
At fixed covariant \(K\), contravariant \(X\), and \(u\), use
\[
 \begin{split}
 R'&=-\langle\Ric,H\rangle+\nabla^i\nabla^jH_{ij}
                                  -\Delta\tr H,\\
 \tau_K'&=-K^{ij}H_{ij},\qquad
 (|K|^2)'=-2(K^2)^{ij}H_{ij},\\
 (P_K^{ij})'&=-H^i{}_aK^{aj}-H^j{}_aK^{ia}
                    +(K^{ab}H_{ab})g^{ij}+\tau_KH^{ij},\\
 (\mathcal L_Xg)'&=\mathcal L_XH,\qquad
 (dV_g)'=\tfrac12(\tr H)dV_g.
 \end{split}
 \]
Integration of \(\mathcal L_XH\) and collection of coefficients gives
\begin{align}
 0={}&\Hess u-(\Delta u)g-u\Ric+2u(K^2-\tau_KK)
      +\frac u2(R-2\Lambda_{\rm c}+\tau_K^2-|K|^2)g \notag\\
 &+\mathcal L_XK-(\operatorname{div}X)K
       -\big(X(\tau_K)+K^{ab}\nabla_aX_b\big)g
       -8\pi X_{(i}J_{j)}.
 \label{eq-unsimplified-metric}
\end{align}
The shift equation gives
\[
 \operatorname{div}X=-u\tau_K,\quad
 K^{ab}\nabla_aX_b=-u|K|^2,\quad
 \tr(\mathcal L_XK)=X(\tau_K)-2u|K|^2.
\]
Taking the trace of \eqref{eq-unsimplified-metric} and using
complementarity yields
\begin{equation}
 \Delta u+X(\tau_K)
   =\frac u2(R+\tau_K^2+|K|^2-4\Lambda_{\rm c}).
 \label{eq-trace-identity}
\end{equation}
Substitution gives \eqref{eq-general-full-lapse}.

To exhibit finite type, write \(B_{ij}=\nabla_{(i}X_{j)}=-uK_{ij}\).
Commuting covariant derivatives gives
\[
 \begin{split}
 \nabla_i\nabla_jX_k
 ={}&\nabla_iB_{jk}+\nabla_jB_{ik}-\nabla_kB_{ij}\\
 &+\tfrac12\big([\nabla_i,\nabla_j]X_k
       -[\nabla_i,\nabla_k]X_j-[\nabla_j,\nabla_k]X_i\big).
 \end{split}
\]
All commutators are curvature times \(X\).  Derivatives of \(B\)
use only \(u,du\) and the first derivatives of \(K\).
Equation~\eqref{eq-general-full-lapse} similarly expresses the
Hessian of \(u\) in first jets.  Therefore, for
\(\mathcal J=(u,X,\nabla u,\nabla X)\),
\begin{equation}
 \nabla\mathcal J=\mathcal A(g,K,\Lambda_{\rm c})\,\mathcal J
 \label{eq-finite-type}
\end{equation}
for a smooth bundle-valued linear coefficient.  ODE uniqueness
along paths proves continuation.  At a zero of \(u\geq0\),
\(du=0\); domination \(|X|\leq u\) also gives \(X=dX=0\).
Nonzero moments consequently have positive lapse everywhere.
\end{proof}

\begin{lemma}[Compactness of normalized adjoint jets]
\label{eq-jet-compactness}
Suppose smooth coefficients in Lemma~\ref{eq-local-adjoint} converge
on compact subsets of a connected open manifold, and nonzero causal
adjoints are defined on an exhaustion containing a fixed point \(p\).
Normalize their full first jets at \(p\) to have norm one.
A subsequence converges smoothly on compact sets to a nonzero causal
adjoint for the limiting coefficients.  The same conclusion holds
for a parameter family on a fixed open manifold.
\end{lemma}

\begin{proof}
The jet at \(p\) cannot be zero by continuation.  On any fixed compact
set, a finite family of coordinate neighborhoods and paths from \(p\)
has bounded lengths and uniformly bounded coefficients in
\eqref{eq-finite-type}.  Grönwall's Inequality gives uniform first-jet
bounds there.  The equations and local elliptic estimates give bounds
of every order on smaller compact sets.  Diagonal compactness gives
the limit.  Its jet at \(p\) still has norm one, and causality is a
closed pointwise condition.  Lemma~\ref{eq-local-adjoint} supplies
positive limiting lapse.
\end{proof}

\subsection{Boundary atoms and hyperbolic normalization}

By Lemma~\ref{eq-area-concentration}, the area term in
\eqref{eq-multiplier-identity} is \(c\,a'_S(H)\).
Thus compact interior tests have zero pairing, and
Lemma~\ref{eq-local-adjoint}, with \(\Lambda_{\rm c}=-3\), applies.
It gives \eqref{eq-shift}--\eqref{eq-full-lapse} in the interior.

\begin{lemma}[One-sided continuation and normalization]
\label{eq-normalization-lemma}
The moment fields extend smoothly to \(S\), their original measure
moments have no part supported on \(S\), and they satisfy
\eqref{eq-end-normalization}.  In particular they are nonzero and
their lapse is positive in the interior.
\end{lemma}

\begin{proof}
In a smooth normal collar of \(S\), the coefficient in
\eqref{eq-finite-type} is smooth up to the face.  Solve its linear
ODE down each collar segment from a fixed interior surface.
Smooth dependence on its footpoint gives a smooth extension; ODE
uniqueness identifies it with the original field in the collar.

Subtract integration by parts of these smooth interior fields from
the compact multiplier identity. Take a metric variation whose value
and full first jet vanish on \(S\), and set the tensor variation to zero.
All area, expansion and integrated boundary terms vanish.
The remaining boundary scalar moment tests
\(-\partial_\nu^2\tr_SH\), which is arbitrary, by choosing a
prescribed \(s^2\) tangential trace in collar coordinates.
Momentum and frame terms involve only lower jets and are zero.
Thus the scalar moment has no boundary part.  Its domination of
the vector moment eliminates that part as well.

We next derive end behavior without imposing growth on the
measures.  Express the first-jet system in a \(b\)-parallel radial
frame, with \(\eta=\operatorname{arsinh}r\).
At \((b,0)\) its lapse system is \(\Hess_bu=ub\) and its shift
system is hyperbolic Killing transport.  Their radial propagators
have norm at most \(C e^{|\eta-\eta'|}\).  This follows directly
from the \(2\)-by-\(2\) blocks \(y''=y\) and the constant blocks,
or from the ambient linear functions and Lorentz generators
on the hyperboloid.
The actual coefficient differs by \(O(e^{-\tau\eta})\);
only two metric derivatives and one tensor derivative occur.
Variation of constants gives
\[
 e^{-\eta}|\mathcal J(\eta)|
 \leq C+C\int_{\eta_0}^{\eta}
        e^{-\tau s}e^{-s}|\mathcal J(s)|\,ds.
 \]
Grönwall therefore gives \(\mathcal J=O(r)\), uniformly in angle.

The lapse and tangential-shift radial equations now read
\[
 y''-y=O(r^{1-\tau}),
\]
and the normal-shift derivative is \(O(r^{1-\tau})\).
Variation of constants, with
\(3/2<\beta<\min\{2,\tau\}\), consequently gives continuous angular
coefficients
\begin{align*}
 u&=r f_0(\omega)+O(r^{1-\beta}),&
 \partial_{\eta}u&=r f_0(\omega)+O(r^{1-\beta}),\\
 X_T&=r W_0(\omega)+O(r^{1-\beta}),&
 X^{\nu_b}&=a_0(\omega)+O(r^{1-\beta}).
\end{align*}
Convergence of the leading coefficients is uniform; no angular
differentiability is required at this stage.  The homogeneous
\(r^{-1}\) mode is included in the remainder because \(\beta<2\).

Use first a conformal prototype \(\psi=r^{-3}f(\omega)\) in the
multiplier identity and integrate to a large sphere.  Its interior
pairing vanishes by the adjoint equations, its \(S\) terms vanish,
and its end-point coefficient is zero.  All pairings converge:
the moments grow at most as \(r\), while the prototype derivatives
are \(O(r^{-4})+O(r^{-3-\tau})\) or better.  The sole nonvanishing
boundary expression is
\[
 -8\int_{r=R}
       (u\partial_{\nu_b}\psi-\psi\partial_{\nu_b}u)\,dA_b.
\]
Terms containing \(KX\psi\) have integrated order \(O(R^{-\tau})\);
replacing \(g\) by \(b\) has the same vanishing order.
The displayed limit is \(32\int_{\mathbb S^2}f_0f\,d\omega\).
The left side is \(16\pi p'_0=32\int_{\mathbb S^2}f\,d\omega\)
by the defining metric flux.  Arbitrary \(f\) gives \(f_0=1\).

For a mixed tensor prototype, the only boundary term is twice
its trace-reversed stress flux against \(X\).  Its limit is
\(2\int_{\mathbb S^2}\langle W_0,W\rangle_\sigma\,d\omega\),
whereas its metric mass and area derivatives vanish.  Hence
\(W_0=0\).  The tangential trace of the shift equation on
coordinate spheres now gives
\[
 \operatorname{div}_{S_R}X_T+H_{S_R}X^{\nu_b}
                      =O(R^{1-\beta}).
\]
Interpret this on the unit sphere.  The divergence term tends to
zero distributionally, since integration against a fixed test
function gains the factor \(R^{-1}\), and \(H_{S_R}\to2\).
Thus \(a_0=0\).  The already uniform limits identify it pointwise.

We have obtained the value estimates in
\eqref{eq-end-normalization}, since \(V_0-r=O(r^{-1})\).
The coupled elliptic system for \((u-V_0,X)\) has uniformly
controlled unit-patch coefficients and source
\(O_{0,\alpha'}(r^{1-\tau})\).  Local estimates applied to these
value bounds yield the first and second derivative Hölder bounds
in \eqref{eq-end-normalization}, with a smaller \(\alpha'\) if
necessary.  The prototype normalization makes the field nonzero;
the last assertion of Lemma~\ref{eq-local-adjoint} gives \(u>0\).
\end{proof}

\subsection{Stationary curvature and free boundary variations}

\begin{proof}[Completion of the proof of Theorem~\ref{eq-causal-killing}]
Since \(u>0\), \eqref{eq-stationary-metric} is a smooth Lorentzian
metric on the open exterior product.  Its coefficients are independent
of \(t\), so \(\xi=\partial_t\) is Killing.  Its future normal is
\(n=u^{-1}(\partial_t-X)\); the future-normal convention fixed after
Equation~\eqref{eq:setup-expansion} gives
\[
 \tfrac12\mathcal L_n g=-\frac1{2u}\mathcal L_Xg=K
\]
by \eqref{eq-shift}.  Thus both induced tensors are the original ones.

We verify its curvature without a spacetime existence assumption.
The Gauss--Codazzi and lapse identities, obtained by differentiating
the displayed metric, give for spatial slots
\[
 \Ric_{\mathbf g,ij}
 =\Ric_{g,ij}+\tau_KK_{ij}-2(K^2)_{ij}
       -u^{-1}\big((\mathcal L_XK)_{ij}+(\Hess u)_{ij}\big),
\]
and
\[
 R_{\mathbf g}=R_g+\tau_K^2+|K|^2
                     -2u^{-1}(\Delta u+X(\tau_K)).
\]
Equation~\eqref{eq-trace-identity} with
\(\Lambda_{\rm c}=-3\) gives \(R_{\mathbf g}=-12\).
Equation~\eqref{eq-full-lapse} then gives
\[
 u(G_{\mathbf g}-3\mathbf g)_{ij}=-8\pi X_{(i}J_{j)}.
\]
The normal and mixed components of this tensor are the constraints,
\(8\pi\mu\) and \(8\pi J_i\).
Complementarity and causality imply
\[
 0=u\mu+J(X)\geq u\mu-|J||X|
                \geq u(\mu-|J|)\geq0.
\]
Where \(\mu>0\), equality forces
\(|X|=u\), \(|J|=\mu\), and \(J=-\mu X^\flat/u\).
Where \(\mu=0\), \(J=0\).
Substitution into all the normal, mixed and spatial components gives
\[
 G_{\mathbf g}-3\mathbf g
      =8\pi\mu u^{-2}\xi^\flat\otimes\xi^\flat,\qquad
 \mu(u^2-|X|^2)=0.
\]
Since the matter term is tracefree, this is exactly
\eqref{eq-null-matter}.  It allows nonzero matter on the null set;
its exclusion is a separate argument in the next section.

It remains to determine the boundary data.  There are no boundary
constraint atoms by Lemma~\ref{eq-normalization-lemma}, and the
integration normal at \(S\) is \(-\nu\).  A compact pure tensor test
\(P\), arbitrary at the boundary, has integrated boundary term
\[
 -2\int_S\big(P(X,\nu)-X^\nu\tr P\big)\,dA_g.
\]
Its expansion derivative is \(\tr_SP\).
Separating mixed and tangential-trace components in
\eqref{eq-multiplier-identity} gives
\[
 X_T=0,\qquad d\eta=2X^\nu\,dA_g.
\]
For a compact simultaneous conformal test
\((4\psi g,2\psi K)\), its area derivative is
\(4\int_S\psi\,dA_g\) and its expansion derivative is
\(4\partial_\nu\psi\).  Integrating the lapse equation gives
\[
 -4c\int_S\psi\,dA_g
 =8\int_S\big[u\partial_\nu\psi
        -(\partial_\nu u+K(X,\nu))\psi\big]\,dA_g
       -4\int_S\partial_\nu\psi\,d\eta .
\]
The boundary value and normal derivative of \(\psi\) can be
prescribed independently.  They give
\(d\eta=2u\,dA_g\), \(X=u\nu\), and
\(\partial_\nu u+K(X,\nu)=c/2=\kappa_h\).

Finally let \(L_{\rm MOTS}\) be the normal expansion linearization
at \(S\), with principal part \(-\Delta_S\).  Extend any
\(s\nu\), \(s\in C^\infty(S)\), to a compact one-sided field \(Y\)
and test with \((\mathcal L_Yg,\mathcal L_YK)\).
In the open exterior its active-ray derivative is zero: under
diffeomorphism transport it differentiates a nonnegative scalar at
a zero, and the difference from isometric transport is annihilated
at activity.  No boundary constraint measure remains.
The mass derivative is zero, while the area and expansion derivatives
are \(\int_SHs\,dA_g\) and \(L_{\rm MOTS}s\).
The multiplier identity therefore gives
\[
 c\int_SHs\,dA_g=2\int_SuL_{\rm MOTS}s\,dA_g,\qquad
 2L_{\rm MOTS}^*u=cH.
\]
Outer area minimization gives \(H\geq0\), by the first area
variation for arbitrary nonnegative outward speeds.
The nonnegative function \(u|_S\) is consequently a supersolution
of an elliptic operator with principal part \(-\Delta_S\).
The strong minimum principle applies: after reversing the operator
and subtracting a sufficiently large zeroth-order constant, the
usual nonpositive zeroth-order condition holds because \(u\geq0\).
On connected \(S\), either \(u>0\) everywhere or \(u=0\) everywhere.
This completes all assertions.
\end{proof}
